\begin{minipage}[t]{0.48\textwidth}

\includegraphics[width=0.65\textwidth]{../../../shared/assets/logo.png}

\end{minipage}
\hfill
\begin{minipage}[t]{0.48\textwidth}

\vspace{-2.0cm}

\raggedleft

{\fontsize{26}{28}\selectfont
\bfseries
\color{HypeOrange}
\edicaoMes
}

\vspace{0.1cm}

{\fontsize{22}{24}\selectfont
\bfseries
\color{HypePurple}
\edicaoAno
}

\end{minipage}

\vspace{0.3cm}

%\begin{center}

% \includegraphics[width=0.5\textwidth]{../shared/assets/capa.png}

% \end{center}


\vspace{1.5cm}

\summaryitem
{01}
{Hype Explica: Erro tipo I vs Erro tipo II}
{Um teste estatístico pode chegar à conclusão errada de duas formas, apontando um efeito que não existe ou deixando passar um efeito real. Entenda a diferença entre os Erros Tipo I e Tipo II e por que equilibrar esses dois riscos é essencial para interpretar resultados.}
{02}

\summaryitem
{02}
{O Papel dos Modelos de Mundo na Próxima Era da IA Física}
{Como a IA pode aprender a imaginar o mundo antes de agir nele? Modelos de mundo permitem que sistemas antecipem consequências, lidem com imprevistos e tomem decisões mais complexas, aproximando a IA de uma atuação mais autônoma no mundo físico.}
{03}

\summaryitem
{03}
{GPT-6 Astra e o Nível
Crítico: Perigos e
Esperanças da IA na
Cibersegurança}
{O GPT-6 Astra chega com capacidades inéditas de atuação autônoma e, segundo a OpenAI, alcança o nível “Crítico” em capacidade cibernética. Com isso, surgem novos desafios para controlar sistemas cada vez mais capazes e, ao mesmo tempo, aproveitar seu potencial para a defesa.}
{04}

\pagebreak

\vspace{0.3cm}